%---------------------------------------------------------------------
% Quick-reference sheet and the recurring mistakes.
%---------------------------------------------------------------------

\section*{فوری دہرائی}
\markboth{فوری دہرائی}{}

\begin{nukta}[آیات کا فوری نقشہ]
\textbf{البقرہ:} \textenglish{30} خلافتِ آدم \ \textbullet\ \
\textenglish{11--13} منافقین \ \textbullet\ \ \textenglish{146,174}
اہلِ کتاب \ \textbullet\ \ \textenglish{112,256} اسلام و آزادیِ دین \
\textbullet\ \ \textenglish{43--44} نماز و زکوٰۃ \ \textbullet\ \
\textenglish{197} حج \ \textbullet\ \ \textenglish{228--229} طلاق و خلع \
\textbullet\ \ \textenglish{224--225} قسمیں \ \textbullet\ \
\textenglish{153--157} صبر و شہادت \ \textbullet\ \
\textenglish{261,275--279} انفاق و سود۔

\smallskip
\textbf{آل عمران:} \textenglish{7} محکمات و متشابہات \ \textbullet\ \
\textenglish{103} اعتصموا بحبل اللہ \ \textbullet\ \ \textenglish{104} امر
بالمعروف \ \textbullet\ \ \textenglish{121--179} غزوہ احد \ \textbullet\ \
\textenglish{78,180} تحریف و بخل \ \textbullet\ \ \textenglish{130} سود \
\textbullet\ \ \textenglish{190--191} اولو الالباب۔
\end{nukta}

\medskip
\subsection*{تین اہم فرق}

\begin{center}
\begin{tabular}{@{}p{3.6cm} p{10.8cm}@{}}
\toprule
\textbf{کافر / منافق} & کھلا انکار / زبان سے اقرار، دل میں انکار --- انجام
  جہنم کا سب سے نچلا درجہ \\
\textbf{لغو / قصدی قسم} & بلا ارادہ نکل جائے، معاف / پختہ ارادے سے کھائی
  جائے، کفارہ لازم \\
\textbf{محکم / متشابہ} & واضح المعنی، عقیدے کی بنیاد / حقیقی مفہوم اللہ کے
  سوا کوئی نہیں جانتا \\
\bottomrule
\end{tabular}
\end{center}

\medskip
\subsection*{چار عام غلطیاں}

\begin{enumerate}\setlength{\itemsep}{3pt}
  \item \textbf{آیت بغیر حوالے کے لکھنا۔} سورت اور آیت کا نمبر ہی سب سے سستا
        نمبر ہے --- کبھی نہ چھوڑیے۔
  \item \textbf{عربی غلط لکھنا۔} جتنی یاد ہے اتنی \emph{صحیح} لکھیے --- آدھی
        صحیح آیت پوری غلط آیت سے بہتر ہے۔
  \item \textbf{کامیابی کے نمبر چالیس سمجھنا۔} ایسوسی ایٹ ڈگری اور بی ایس کے
        پرچوں میں یہ \textbf{پچاس} ہیں۔
  \item \textbf{تفسیر کے بغیر صرف ترجمہ لکھنا۔} سوال ''ترجمہ و تفسیر'' مانگتا
        ہے --- ترجمے کے بعد دو تین سطریں تشریح کی لازمی لکھیے۔
\end{enumerate}

\medskip
\noindent
\textbf{\color{bmblue}اگر وقت بہت کم ہو۔}\ \
وہ موضوعات جو \textbf{تین یا چار} پرچوں میں آئے، پہلے تیار کیجیے: اہلِ کتاب
کی خصوصیات و تحریف (سوال \textenglish{3}، \textenglish{14})، سود اور قسموں
کے احکام (سوال \textenglish{8}، \textenglish{17})، طلاق و خلع (سوال
\textenglish{7})، اور الفاظ کی پہچان (سوال \textenglish{18})۔ ماخذ کی تفصیل
کتابچے کے آغاز میں درج ہے۔
